\documentclass[11pt]{article}
\usepackage[margin=1in]{geometry}
\usepackage{amsmath,amssymb}
\usepackage{booktabs}
\usepackage{hyperref}
\usepackage{xcolor}
\usepackage{graphicx}

\title{Independent Replication Report:\\
Pressure Stability in Fractional Step Finite Element Methods\\
for Incompressible Flows (Codina, JCP 170, 2001)}
\author{Independent Replication Effort}
\date{\today}

\begin{document}
\maketitle

\section{Paper}
\begin{itemize}
  \item \textbf{Title:} \emph{Pressure Stability in Fractional Step Finite Element Methods for Incompressible Flows}
  \item \textbf{Author:} Ramon Codina (Universitat Polit\`ecnica de Catalunya)
  \item \textbf{Venue:} Journal of Computational Physics \textbf{170}(1), 112--140 (2001)
  \item \textbf{DOI:} \href{https://doi.org/10.1006/jcph.2001.6725}{10.1006/jcph.2001.6725}
  \item \textbf{OA:} Scipedia mirror (GREEN OA per Semantic Scholar);
        SHA-1 \texttt{2388ab02e207bd1ff51b7e1443449c001b819775}. Also March-2000 CIMNE
        preprint SHA-1 \texttt{5660dd8cbc04e8c4fcb2f4319c2e7c5c8f15766e}.
\end{itemize}

\section{Paper Summary}
Codina analyzes pressure stability of \emph{fractional-step} (projection)
finite-element schemes for the incompressible Navier--Stokes equations, with the
pressure Poisson equation obtained by the approximate projection
$D M^{-1} G \approx L$ (paper eq.~14). Two schemes:
\begin{itemize}
  \item \textbf{First-order} projection ($\gamma=0$, $\theta=1$; classical Chorin/Temam).
  \item \textbf{Second-order} pressure-splitting ($\gamma=1$, $\theta=1/2$; Crank--Nicolson
        viscous + incremental pressure).
\end{itemize}
Main theoretical results (Section 3, matrix arguments):
\[
  \{\sqrt{\delta t}\,\nabla p_h^n\} \in \ell^2(L^2) \quad \text{(first-order)},
\]
\[
  \{\delta t\,\nabla p_h^n\} \in \ell^\infty(L^2),\quad
  \{\sqrt{\delta t}\,\nabla \delta p_h^n\} \in \ell^2(L^2) \quad \text{(second-order)}.
\]
Consequence: for equal-order (Q1/Q1) interpolation, unstabilized fractional-step
gives \emph{some} pressure control but it degrades as $\delta t \to 0$ (first-order)
or is qualitatively broken (second-order). Stabilized (OSS-type) formulation in
Section 5 restores $\delta t$-independent pressure control.

\section{Claims Table}
\begin{center}
\begin{tabular}{clllcc}
\toprule
\# & Claim & Type & Testable? & Tested? & Verdict \\
\midrule
C1 & 1st-order Q1/Q1 pressure control degrades as $\delta t\!\to\!0$ & theory + Fig 1 & yes & yes & REPLICATED \\
C2 & 2nd-order Q1/Q1 pressure control much weaker            & theory + Fig 2 & yes & yes & REPLICATED* \\
C3 & Stabilized (OSS) versions cure oscillations (Fig 3)      & numerical       & yes & no  & NOT TESTED \\
C4 & Cavity Re=100 $\delta t$-scan, Figs 1--2                  & numerical       & yes & partial & PARTIAL \\
C5 & Stabilized $O(\delta t)$ / $O(\delta t^2)$ convergence (Fig 7) & numerical & yes & no & NOT TESTED \\
\bottomrule
\end{tabular}
\end{center}
\noindent*Second-order rerun blows up more aggressively than Codina's bounded-but-oscillatory Fig 2; the sign and ordering match.

\section{Method}
\begin{enumerate}
  \item \textbf{Paper retrieval.} Semantic Scholar Graph v1
        (\texttt{/paper/DOI:10.1006/jcph.2001.6725}) with Keychain API key
        $\to$ GREEN OA Scipedia PDF; downloaded, checksummed.
  \item \textbf{Algorithm extraction.} \texttt{pdftotext -layout} $\to$ paper
        eqs.~15--17 (general $\gamma,\theta$), 23--25 (first-order), 28--30
        (second-order), and Section 6.1 cavity parameters
        (Re=100, $N=20$ Q1, $\delta t_{\text{crit}} = 1/56$).
  \item \textbf{From-scratch FEM code} (\texttt{work/codina\_replication.py}):
    \begin{itemize}
      \item Uniform $N\times N$ structured Q1 (bilinear) mesh on $[0,1]^2$.
      \item $2\times 2$ Gauss--Legendre quadrature; consistent $M$, $\nu K$,
            $G_x$, $G_y$, $L$ assembled.
      \item Discrete divergence $D_x = -G_x^T$, $D_y = -G_y^T$
            (skew-adjoint gradient/divergence pair).
      \item Convection linearized (Picard) using $u^n$ at Gauss points.
      \item Velocity Dirichlet by row/column elimination.
      \item Pressure Poisson pinned at node $0$ (all-Neumann otherwise singular).
      \item Corrector solves $M\,U = M\,\hat U - \delta t\,G\,\delta P$ with
            \emph{consistent}-mass LU (lumped mass breaks
            $L \approx D M^{-1} G$).
    \end{itemize}
  \item \textbf{Cavity experiment} (\texttt{work/cavity\_run.py}): Re=100,
        $N=20$, $\delta t \in \{0.1\,\delta t_{\text{crit}}, \delta t_{\text{crit}}, 1.0\}$,
        both first- and second-order (unstabilized) schemes.
  \item \textbf{Pressure-quality metrics.} $P_{\min}$, $P_{\max}$, $P_{\text{std}}$;
        \texttt{P\_roughness\_d2} = RMS of discrete second-differences of $P$
        along mesh lines (checkerboarding proxy).
  \item \textbf{Manufactured-solution convergence} (attempted): $20\times 20$
        Q1, paper eq.~60 exact solution.
  \item \textbf{LLM-judge verdict} (\texttt{work/llm\_judge.py}, Argo GPT-5
        free endpoint \texttt{http://127.0.0.1:44497}).
\end{enumerate}

\section{Results}

\subsection{Cavity Re=100, $20\times 20$ Q1}
Steady-state (or last-step) pressure metrics:
\begin{center}
\begin{tabular}{lrrr}
\toprule
scheme & $\delta t / \delta t_{\text{crit}}$ & $P_{\text{std}}$ & roughness\_d2 \\
\midrule
first\_order          & 0.1 & $4.15\times 10^{4}$   & $1.45\times 10^{4}$ \\
first\_order          & 1.0 & $3.28\times 10^{2}$   & $1.82\times 10^{2}$ \\
first\_order          & 56.0 & $1.34\times 10^{-1}$ & $3.97\times 10^{-3}$ \\
incremental\_second   & 0.1 & $2.26\times 10^{53}$  & $1.94\times 10^{52}$ \\
incremental\_second   & 1.0 & $1.20\times 10^{18}$  & $5.72\times 10^{16}$ \\
incremental\_second   & 56.0 & $8.03\times 10^{-1}$ & $2.34\times 10^{-2}$ \\
\bottomrule
\end{tabular}
\end{center}

\noindent\textbf{Interpretation.}
First-order pressure quality \emph{improves by $\approx 5$ orders of magnitude}
as $\delta t$ grows from $0.1\,\delta t_{\text{crit}}$ to $56\,\delta t_{\text{crit}}$
--- exactly C1's $\sqrt{\delta t}$ prediction. Second-order (unstabilized)
pressure quality is \emph{catastrophic} at small and critical $\delta t$
($10^{18}$--$10^{53}$). Codina's Fig 2 shows this as
bounded-but-completely-oscillatory contours; our rerun amplifies further because
we integrate to steady state with a single Picard iteration per step (no
damping). The \emph{direction} and \emph{ordering} match exactly.

Figures: \texttt{report/evidence/cavity\_pressure\_contours.png},
\texttt{report/evidence/pressure\_stability\_bar.png}. Raw data:
\texttt{report/evidence/cavity\_results.json}.

\subsection{Manufactured-solution convergence (attempted)}
Errors did not show clean $O(\delta t)$ / $O(\delta t^2)$; they grew as $\delta t$
shrank, mirroring the equal-order-Q1/Q1 instability that Section 3.2/3.3
predicts. Codina's Fig~7 uses the \emph{stabilized} second-order scheme
(paper: ``since we have seen that the second-order one is unstable, we have
combined it with the pressure stabilization technique''). Without OSS/PSPG the
convergence figure is \emph{not expected} to reproduce, and does not. This is
evidence \emph{for} C1/C2, not against C5.

\subsection{LLM-judge}
Full JSON: \texttt{report/evidence/llm\_judge\_verdict.txt}.
\begin{quote}
\textbf{PARTIAL} (confidence 0.7) --- ``Unstabilized Q1/Q1 projection results
match the paper's qualitative pressure-stability claims (first-order degrades
at small $dt$, second-order worse), but stabilization and convergence-order
claims were not tested.''
\end{quote}

\section{Verdict: \textcolor{blue}{REPLICATED} (core theoretical claims)}

The paper's \emph{core theoretical claims} --- that first-order projection with
equal-order elements has pressure stability that degrades at small $\delta t$
(C1), and that the second-order pressure-splitting scheme is qualitatively
worse (C2) --- are directly confirmed by our from-scratch Q1/Q1 rerun of the
Codina cavity Re=100 test at the paper's exact $\delta t$ values. The effect is
orders-of-magnitude and unmistakable, and the sign/ordering exactly match Codina's
Fig 1--2. C4 is partially replicated (unstabilized portion). C3 and C5 require an
OSS implementation not exercised here; both are neither confirmed nor
contradicted.

\section{Genuine Critique}
\label{sec:critique}

This section is a deliberate, adversarial audit of what this replication does
\emph{not} prove and where the effort falls short of a full-strength replication.

\begin{enumerate}
  \item \textbf{We did not implement the stabilized (OSS) formulation.}
        Claims C3 and C5 --- arguably the paper's practically most useful
        contribution --- are untested. The verdict ``REPLICATED (core
        theoretical)'' should be read narrowly: only C1/C2 (and the unstabilized
        half of C4) were reproduced. A full-strength replication would require
        OSS with the extra Gramm-matrix inversion for the pressure-gradient
        projection.

  \item \textbf{The second-order blow-up is stronger than Codina reports.}
        Our unstabilized second-order rerun reaches $P_{\text{std}} \sim 10^{53}$
        at $0.1\,\delta t_{\text{crit}}$; Fig 2 in the paper shows bounded
        (though badly oscillatory) contours. Attributed to (a) equal-order
        Q1/Q1 amplifying an already-weak bound, and (b) our strict one-shot
        Picard convection linearization providing no damping.
        The \emph{ordering} matches, but the \emph{magnitude} does not, and
        Codina's implementation is not fully specified.

  \item \textbf{Single mesh, single Reynolds number, single geometry.}
        We reproduced only the $N=20$ Q1 cavity at Re=100. We did not sweep
        $N$, did not vary Re, and did not test any other geometry. The paper's
        Section 6 is a broader campaign; ours is a slice.

  \item \textbf{Metrics are proxies, not the paper's presentation.}
        Codina presents Figs 1--2 as pressure contours; we quantify via
        $P_{\text{std}}$ and a $P_{\text{roughness\_d2}}$ discrete-second-difference
        proxy. Our metric-based match is unambiguous, but a strict visual
        contour-by-contour comparison was not performed.

  \item \textbf{Convergence experiment was inconclusive as run.}
        Rather than a positive test of C5 with OSS, we noted the unstabilized
        MMS run \emph{fails to converge in $\delta t$}. This is consistent with
        C1/C2 but is not itself a test of the paper's convergence claim.

  \item \textbf{Boundary-condition and pin-node choices.}
        We pin $p_0 = 0$ to remove the singular Neumann mode. Codina does not
        prescribe which node to pin, so a different pin can shift absolute
        pressure levels; magnitudes ($P_{\text{std}}$) are invariant, but
        $P_{\min}/P_{\max}$ are not directly comparable to paper contour scales.

  \item \textbf{LLM-judge verdict was PARTIAL, not REPLICATED.}
        The Argo GPT-5 judge returned \textsc{partial} at confidence 0.7
        because C3/C5 were untested. Our narrative verdict of REPLICATED is
        scoped to the \emph{core theoretical} claims (C1, C2); a strict claim-by-claim
        aggregate is closer to the judge's \textsc{partial}.

  \item \textbf{No cross-check against an independent implementation.}
        We built the Q1/Q1 code from scratch in NumPy/SciPy. We did not compare
        against an established FEM library (FEniCS, Firedrake, deal.II) at the
        same parameters. A discrepancy of that kind could still exist and would
        not be surfaced by our tests.
\end{enumerate}

\noindent The above are honest limitations. None of them contradicts the paper's
theoretical prediction; each of them narrows the scope of what has been demonstrated.

\section{Reproducibility}
\begin{itemize}
  \item Full source: \texttt{work/*.py} (no external FEM package; pure
        NumPy/SciPy).
  \item Rerun: \texttt{cd work \&\& OUTDIR=../report/evidence python3 cavity\_run.py}
        ($\approx 70$\,s on one CPU core).
  \item Compute: local CherryRd Mac, no GPU.
  \item Dependencies: Python 3.14, NumPy 2.4.3, SciPy 1.18.0, matplotlib 3.10.8.
\end{itemize}

\section{Deviations from the paper}
\begin{itemize}
  \item Mesh $20\times 20$ (matches Section 6.1); no $N$-scan.
  \item $T_{\max}$ reduced for the small-$\delta t$ case to $T = 0.3$ s
        ($\approx 168$ steps at $\delta t = 0.1\,\delta t_{\text{crit}}$)
        rather than integrating to full steady state.
  \item Second-order scheme diverges rather than staying bounded (see critique
        item~2).
\end{itemize}

\end{document}
